\documentclass{article}
\usepackage{amsmath}

\begin{document}

Maxwell's Equations are a set of four fundamental equations in electromagnetism that describe how electric and magnetic fields interact. In their differential form, they are expressed as:

1. \textbf{Gauss's Law for Electricity}:
\[
\nabla \cdot \mathbf{E} = \frac{\rho}{\varepsilon_0}
\]
This equation states that the electric flux through a closed surface is proportional to the charge enclosed.

2. \textbf{Gauss's Law for Magnetism}:
\[
\nabla \cdot \mathbf{B} = 0
\]
This indicates that there are no magnetic monopoles; magnetic field lines are continuous and have no beginning or end.

3. \textbf{Faraday's Law of Induction}:
\[
\nabla \times \mathbf{E} = -\frac{\partial \mathbf{B}}{\partial t}
\]
This describes how a time-varying magnetic field induces an electric field.

4. \textbf{Ampère-Maxwell Law}:
\[
\nabla \times \mathbf{B} = \mu_0 \mathbf{J} + \mu_0 \varepsilon_0 \frac{\partial \mathbf{E}}{\partial t}
\]
This relates magnetic fields to the currents and changing electric fields that produce them.

Where:
\begin{itemize}
  \item $\mathbf{E}$ is the electric field,
  \item $\mathbf{B}$ is the magnetic field,
  \item $\rho$ is the electric charge density,
  \item $\mathbf{J}$ is the current density,
  \item $\varepsilon_0$ is the permittivity of free space,
  \item $\mu_0$ is the permeability of free space.
\end{itemize}

\end{document}